\section{Análise dos Resultados}
\label{sec:analysis}

\subsection{Tempo, escalabilidade e estrutura}

\textbf{Resultado observado.}
A aproximação tem tempo menor nas \FasterApproxInstances{} de
\TotalInstances{} comparações de execuções selecionadas, incluindo
os grafos sem arestas. No exato, as medianas dos grupos aleatórios
crescem com $n$ na Tabela~\ref{tab:sizes}. Para $p=0{,}5$,
a diferença entre os extremos da grade é de várias ordens de
grandeza, enquanto a aproximação permanece na escala de
microssegundos.

\textbf{Explicação possível.}
A aproximação faz uma passagem pelas arestas e usa marcações
por vértice. O exato, além do custo de sua incumbente inicial,
recalcula graus e limites em numerosos estados e copia conjuntos
de bits. A ampliação da árvore de busca observada é coerente
com o comportamento exponencial possível dessa estratégia.
Não se trata de uma demonstração experimental da notação
assintótica: a justificativa de complexidade é algorítmica.

\textbf{Limitação.}
Os tempos aproximados não crescem monotonicamente com $n$.
Por exemplo, a mediana agregada é 4,90~$\mu$s para $n=10$
e 2,20~$\mu$s para $n=30$. Esse resultado não indica que o
algoritmo se torna assintoticamente mais rápido. A coleta usa
uma única JVM, instâncias em ordem fixa e grafos pequenos;
otimização progressiva pelo compilador em tempo de execução,
resolução do relógio, frequência da CPU e coleta de lixo
podem afetar tempos tão curtos. O aquecimento reduz, mas
não elimina esses fatores. Não houve múltiplas inicializações
independentes da JVM nem cálculo de significância estatística.

Na Tabela~\ref{tab:configurations}, aumentar $p$ com $n$ fixado
está associado a maior tempo mediano do exato em todos os
cinco tamanhos desta grade. A Figura~\ref{fig:density}, porém,
não autoriza concluir que densidade sozinha determina a
dificuldade. O bipartido completo com 30 vértices tem
densidade aproximadamente 0,517 e tempo exato de 0,1362~ms,
enquanto os aleatórios de $n=30,p=0{,}5$, de densidade média
0,496, têm mediana de 560,041~ms. O arranjo das arestas e a
qualidade dos limites também importam.

\subsection{Otimalidade, limites e timeout}

Os 228 resultados \texttt{OPTIMAL} indicam prova produzida pelo
solver; os dois \texttt{TIME\_LIMIT} indicam apenas incumbentes
válidas. Essa distinção é necessária mesmo em uma família
com ótimo teórico conhecido. Em $K_{30}$, o tamanho 29
coincide com a fórmula $n-1$, mas o algoritmo interrompido
não concluiu sua própria certificação. O resultado continua
classificado como timeout, e seus campos de ótimo, razão e
gap permanecem \texttt{NA}.

Uma explicação para o custo dos completos está na força do
limite inferior. Em um completo residual com $r$ vértices,
um matching tem $\lfloor r/2\rfloor$ arestas, enquanto o
mínimo residual é $r-1$. O limite pode, portanto, ficar
distante da quantidade realmente necessária, retardando
podas. Já uma primeira cobertura de boa qualidade e limites
mais próximos do mínimo favorecem outros grafos.
A implementação não utiliza reconhecimento de completos,
algoritmos especializados para bipartidos ou memorização
de subproblemas residuais repetidos.

As contagens de nós e podas sustentam a descrição do esforço
da busca, mas não medem isoladamente o benefício de cada
heurística. Não foram executadas versões sem limite
inferior, sem ordenação por grau ou sem incumbente
aproximada. Assim, atribuir uma aceleração quantitativa
a uma dessas técnicas exigiria uma comparação adicional.
Além disso, os nós e tempos de timeout são observações
censuradas: não permitem estimar quantos estados ou quanto
tempo faltava para concluir a prova.

\subsection{Qualidade e compromisso entre os algoritmos}

A aproximação respeita a garantia de fator~2 em todas as
223 instâncias em que a razão pode ser avaliada.
Seu fator médio de \MeanRatio{} corresponde ao gap médio
de \MeanGap{}\%, e \FactorTwoInstances{} instâncias atingem
exatamente o limite de fator~2. Caminhos, estrelas e
bipartidos completos da suíte tornam visível que a garantia
não é apenas um limite folgado sem ocorrência prática
nesses dados.

As coincidências com o ótimo em $K_{15}$ e $K_{25}$
decorrem da cardinalidade dos extremos emparelhados nesses
completos ímpares: o método deixa um vértice livre e
seleciona $n-1$. São resultados particulares, não uma
capacidade geral de reconhecer quando a solução é ótima.
Nos grafos sem arestas, a coincidência é trivial e deve
ser separada da análise de razão.

A composição da amostra também influencia a média:
o subconjunto aleatório elegível tem média
\RandomMeanRatio{}, enquanto o estruturado elegível tem
\StructuredMeanRatio{} (Tabela~\ref{tab:quality}).
Essa diferença decorre das famílias e sementes escolhidas,
e não demonstra superioridade estatística em uma população
de grafos reais. Os dois timeouts não participam das
estatísticas de qualidade; usar suas incumbentes como
denominador produziria uma razão sem a interpretação
de distância ao ótimo requerida.

O Branch and Bound oferece a vantagem do certificado
quando conclui, ao custo de tempo potencialmente elevado.
O matching maximal oferece execução curta e uma garantia
a priori, mas pode selecionar mais vértices. A escolha
adequada depende de se exigir mínimo comprovado ou de
se aceitar uma solução limitada pelo fator~2. A amostra,
restrita a grafos sintéticos de até 30 vértices e a uma
máquina, não sustenta superioridade universal, extrapolação
para redes maiores ou conclusões sobre modelos ponderados.
